% Frozen refinement module; requires hko-ascent-main.tex notation and results.
\subsection{Seventeen facets and the exact response to unequal cut speeds}
\label{subsec:hko-seventeen}

The twenty-row family admits a sharper facet-count conclusion. Keep the ten
profiles of~\eqref{eq:ha-profile} fixed, but give their new rows nonnegative
speeds $w=(w_0,\ldots,w_9)$:
\[
 H_{t,w}=\{x:a_i\cdot x\leq1,\
                 (a_i+tw_iD_i)\cdot x\leq1,\ 0\leq i<10\}.
\]
Its profile on facet $i$ is $w_if_i$. Write
\[
 \kappa=\int_{F_i}f_i\,d\nu
        =\frac{8(5+\sqrt5)}{46875};
\]
this value is independent of the facet and includes its cone-volume mass
$1/10$.

\begin{proposition}[Exact weighted response]\label{prop:ha-weighted}
For every $w_i\geq0$, the supremum over all finite maximizing words and their
weak occupation limits is
\begin{equation}\label{eq:ha-weighted-response}
 \sup_\gamma\sum_{i=0}^9w_i\int_{F_i}f_i\,d\mu_\gamma
       =\frac{\max_iw_i}{125}.
\end{equation}
Consequently
\begin{equation}\label{eq:ha-weighted-derivative}
 D^+\log\operatorname{sys}(H_{t,w})
       =4\left(\kappa\sum_iw_i-\frac{\max_iw_i}{125}\right).
\end{equation}
\end{proposition}

\begin{proof}
Nonnegativity and Proposition~\ref{prop:ha-certificate}, together with
Lemma~\ref{lem:ha-repeat}, give the upper bound
\[
 \sum_iw_i\int_{F_i}f_i\,d\mu_\gamma
 \leq(\max_iw_i)\int f\,d\mu_\gamma
 \leq\frac{\max_iw_i}{125}.
\]
For the reverse bound we exhibit a maximizing word whose entire response
is concentrated on one facet. Put $u=\tau^2$ and take
\begin{align*}
 \sigma&=(0,6,5,2,1,8,4),\\
 \beta&=\frac1{200}
 (17+5u,\ 25+5u,\ 25+5u,\ 44-8u,\ 8,\ 50-10u,\ 31+3u).
\end{align*}
The weights are strictly positive, sum to one, close the normal word, and
give $Q=\tau/5$. Its KKT reconstruction has constant radial hinge arguments,
in this order,
\[
 \ell_{\sigma_r}-\frac35
 =\frac1{40}(-59+5u,\ -9-5u,\ -9-5u,\ -9-5u,\
                       8,\ -9-5u,\ -45+3u).
\]
Since $0<u<1$, only the facet labeled $1$ contributes. Its weight is
$1/25$ and its hinge value is $1/5$, giving response $1/125$ on that facet
and zero on all others. These are literal exact predicates checked by the
supplementary witness replay; the same point is chart~11 of the complete
catalogue, at parameter $3/10-\tau^2/50$.

Simultaneous rotation of the q- and p-planes by $2\pi/5$ preserves $H$,
the symplectic form, and this labeled profile family. The symplectic
involution
\[
 S(q,p)=(J_2p,-J_2q)
\]
also preserves $H$ and exchanges facet $i$ with facet $5+i$, including
their cut velocities. These symmetries act transitively on the ten facets.
Thus every facet admits a shortest word with response $1/125$ concentrated
there. Choose a facet with maximal $w_i$ to obtain the lower bound in
\eqref{eq:ha-weighted-response}. The derivative formula follows from
Theorem~\ref{thm:mc-derivative} and the volume response
$\kappa\sum_iw_i$.
\end{proof}

\begin{corollary}[A seventeen-facet ascent]\label{cor:ha-seventeen}
Choose any seven of the ten cuts in~\eqref{eq:ha-directions}, with unit
speed, and omit the other three new rows. For all sufficiently small
positive $t$ the resulting body has seventeen genuine facets, and
\[
 D^+\log\operatorname{sys}(H_{t,w})
 =\frac{4(56\sqrt5-95)}{46875}>0.
\]
Within the cone of nonnegative speeds on these fixed ten profiles, seven
new rows is the smallest number that can give a positive first variation.
\end{corollary}

\begin{proof}
The exact comparisons are
\[
 6\kappa=\frac{240+48\sqrt5}{46875}
   <\frac{375}{46875}=\frac1{125}
   <\frac{280+56\sqrt5}{46875}=7\kappa.
\]
Seven unit speeds therefore give the displayed positive derivative. Every
selected new row and every old row is a genuine facet by the open-patch
argument following~\eqref{eq:ha-family}, so the facet count is seventeen.
If instead at most six speeds are nonzero and $w\ne0$, then
$\sum_iw_i\leq6\max_iw_i$, making
\eqref{eq:ha-weighted-derivative} strictly negative. This proves the claimed
minimal count for this fixed profile cone; it does not exclude ascent using
different profiles or a different perturbation with fewer new facets.
\end{proof}
